\ifdefined\gkver\else\def\gkver{student}\fi
\documentclass[\gkver]{gaokaozhenti}
\begin{document}

\chapter{2004年上海理综}
%% paperType: 理综物理部分

\begin{enumerate}
\item
%% number: 10
%% paperName: 2004年上海理综第 10 题 3 分
%% typeId: 1
%% type: 单选题
%% chapter: 电磁感应
%% point: 电磁感应现象
%% method: 无
%% score: 3
%% degree: 950
%% duplicateId: 0
%% body:
发电的基本原理之一是电磁感应，发现电磁感应现象的科学家是\xzanswer{C}

\fourchoices[answer=C]
{安培}
{赫兹}
{法拉第}
{麦克斯韦}

%% answer: C
%% memo:
\memoanswer{
本题原卷及材料中未提供详解，待补充。
}

\item
%% number: 11
%% paperName: 2004年上海理综第 11 题 3 分
%% typeId: 1
%% type: 单选题
%% chapter: 电路
%% point: 电功电功率
%% method: 无
%% score: 3
%% degree: 750
%% duplicateId: 0
%% body:
$1\UkWh$ 电能可以做很多事。请估算 $1\UkWh$ 的电能全部用来托起一位普通高中生，使他提升的高度最接近\xzanswer{D}

\fourchoices[answer=D]
{$2\Um$}
{$20\Um$}
{$700\Um$}
{$70000\Um$}

%% answer: D
%% memo:
\memoanswer{
本题原卷及材料中未提供详解，待补充。
}

\item
%% number: 12
%% paperName: 2004年上海理综第 12 题 3 分
%% typeId: 1
%% type: 单选题
%% chapter: 电路
%% point: 电功电功率
%% method: 无
%% score: 3
%% degree: 550
%% duplicateId: 0
%% body:
分时电表惠及百万家。2001 年上海市启用分时电表，家庭用电在谷时段（22:00$\sim$次日 6:00）以每千瓦时 0.30 元优惠计费，平时段（6:00$\sim$22:00）仍以每千瓦时 0.61 元计费。某家庭五月份平时段和谷时段的用电数分别为 $116\UkWh$ 和 $98\UkWh$。分时电表的安装，为该家庭节省人民币\xzanswer{B}

\fourchoices[answer=B]
{29.40 元}
{30.38 元}
{34.80 元}
{35.96 元}

%% answer: B
%% memo:
\memoanswer{
本题原卷及材料中未提供详解，待补充。
}

\item
%% number: 13
%% paperName: 2004年上海理综第 13 题 3 分
%% typeId: 1
%% type: 单选题
%% chapter: 电路
%% point: 电功电功率
%% method: 无
%% score: 3
%% degree: 850
%% duplicateId: 0
%% body:
如今家用电器越来越多，它们在待机状态下也会耗电。为了节约电能和用电安全，你将选择\xzanswer{C}

①切断长期不使用的空调机的电源　②电视机在不使用时，不切断电源　③合理使用节能灯具　④尽量避免同时使用大功率家用电器

\fourchoices[answer=C]
{①②③}
{②③④}
{①③④}
{①②④}

%% answer: C
%% memo:
\memoanswer{
本题原卷及材料中未提供详解，待补充。
}

\item
%% number: 14
%% paperName: 2004年上海理综第 14 题 3 分
%% typeId: 1
%% type: 单选题
%% chapter: 电路
%% point: 电功电功率
%% method: 无
%% score: 3
%% degree: 650
%% duplicateId: 0
%% body:
抽水蓄能电站可以在用电低谷时，将电网中多余的电能转化为重力势能；在用电高峰时，再将重力势能转化为电能输回电网，假定多余电能为 $5\times10^{4}\UkWh$，这部分多余电能经储存再回到电网过程中将损失 44\%，则此蓄能电站能给电网补充的电能为\xzanswer{B}

\fourchoices[answer=B]
{$2.2\times10^{4}\UkWh$}
{$2.8\times10^{4}\UkWh$}
{$3.5\times10^{4}\UkWh$}
{$5\times10^{4}\UkWh$}

%% answer: B
%% memo:
\memoanswer{
本题原卷及材料中未提供详解，待补充。
}

\item
%% number: 34
%% paperName: 2004年上海理综第 34 题 6 分
%% typeId: 3
%% type: 填空题
%% chapter: 直线运动
%% point: 多段匀变速直线运动
%% method: 无
%% score: 6
%% degree: 550
%% duplicateId: 0
%% body:
为了安全，在行驶途中，车与车之间必须保持一定的距离。因为，从驾驶员看见某一情况到采取制动动作的时间里，汽车仍然要通过一段距离（称为思考距离）；而从采取制动动作到车完全停止的时间里，汽车又要通过一段距离（称为制动距离）。下表给出了汽车在不同速度下的思考距离和制动距离等部分数据。请分析这些数据，完成表格。

\begin{center}
\begin{tblr}{|c|c|c|c|}
\hline
速度（$\Ukmh$） & 思考距离（m） & 制动距离（m） & 停车距离（m） \\
\hline
45 & 9 & 14 & 23 \\
\hline
75 & 15 & 38 & \tkanswer[1cm]{53} \\
\hline
90 & \tkanswer[1cm]{18} & \tkanswer[1cm]{55} & 73 \\
\hline
105 & 21 & 75 & 96 \\
\hline
\end{tblr}
\end{center}

%% answer:
%% memo:
\memoanswer{
本题原卷及材料中未提供详解，待补充。
}

\item
%% number: 35
%% paperName: 2004年上海理综第 35 题 3 分
%% typeId: 1
%% type: 单选题
%% chapter: 运动和力
%% point: 牛顿定律的应用
%% method: 无
%% score: 3
%% degree: 650
%% duplicateId: 0
%% body:
在行车过程中，如果车距不够，刹车不及时，汽车将发生碰撞，车里的人可能受到伤害，为了尽可能地减轻碰撞引起的伤害，人们设计了安全带。假定乘客质量为 $70\Ukg$，汽车车速为 $108\Ukmh$（即 $30\Ums$），从踩下刹车到车完全停止需要的时间为 $5\Us$，安全带对乘客的作用力大小约为\xzanswer{A}

\onepicture[w=4.5cm, align=right]{figs/35.jpg}

\fourchoices[answer=A]
{$400\UN$}
{$600\UN$}
{$800\UN$}
{$1000\UN$}

%% answer: A
%% memo:
\memoanswer{
本题原卷及材料中未提供详解，待补充。
}

\end{enumerate}

\end{document}
